\documentclass{article}
\usepackage[utf8]{inputenc}
\usepackage[spanish]{babel}
\usepackage{listings}
\usepackage{xcolor}
\usepackage{geometry}
\geometry{a4paper, margin=1in}

\% Configuración para código C++
\definecolor{codegreen}{rgb}{0,0.6,0}        
\definecolor{codegray}{rgb}{0.5,0.5,0.5}     
\definecolor{codepurple}{rgb}{0.58,0,0.82}   
\definecolor{backcolour}{rgb}{0.95,0.95,0.92}

\lstset{
    backgroundcolor=\color{backcolour},      
    commentstyle=\color{codegreen},
    keywordstyle=\color{magenta},
    numberstyle=\tiny\color{codegray},       
    stringstyle=\color{codepurple},
    basicstyle=\ttfamily\small,
    breakatwhitespace=false,
    breaklines=true,
    captionpos=b,
    keepspaces=true,
    numbers=left,
    numbersep=5pt,
    showspaces=false,
    showstringspaces=false,
    showtabs=false,
    tabsize=2,
    language=C++,
    inputencoding=utf8,
    extendedchars=true,
    literate={á}{{\'a}}1 {é}{{\'e}}1 {í}{{\'i}}1 {ó}{{\'o}}1 {ú}{{\'u}}1 {ñ}{{\~n}}1 {¿}{{\textquestiondown}}1 {¡}{{\textexclamdown}}1
}

\% Macros para las secciones de los apuntes
\newcommand{\quees}[1]{
    \section*{🤔 ¿QUÉ ES?}
    \hrulefill \\
    \#1
}

\newcommand{\dichodeotraforma}[1]{
    \section*{💬 DICHO DE OTRA FORMA}
    \hrulefill \\
    \#1
}

\newcommand{\diagrama}[1]{
    \section*{🎨 DIAGRAMA}
    \hrulefill \\
    \#1
}

\newcommand{\comofunciona}[1]{
    \section*{⚙️ ¿CÓMO FUNCIONA?}
    \hrulefill \\
    \#1
}

\newcommand{\complejidad}[1]{
    \section*{📱 COMPLEJIDAD (Big-O)}
    \hrulefill \\
    \#1
}

\newcommand{\entrevistas}[1]{
    \section*{🎯 EN ENTREVISTAS}
    \hrulefill \\
    \#1
}

\newcommand{\resumen}[1]{
    \section*{📋 RESUMEN RÁPIDO}
    \hrulefill \\
    \#1
}

\begin{document}

\title{Union-Find (Disjoint Set Union - DSU)}
\author{Aldo Zepeda Montoya}
\date{\today}
\maketitle

\subsection*{CATEGORÍA: 02\_Estructuras\_No\_Lineales}

\quees{
Los clubes de la escuela 🏫. Al principio, cada estudiante es su
propio club (un conjunto de 1 persona). 

Cuando dos personas se hacen amigas, sus clubes se UNEN. Para saber 
si dos personas están en el mismo club, sigues la cadena de amistades
hasta llegar al "líder" del grupo. Si ambos tienen el mismo líder,
entonces ¡están en el mismo club!

Es una estructura que gestiona grupos de elementos que no se 
traslapan y te permite unirlos o preguntar si están conectados 
muy rápido.
}

\dichodeotraforma{
Union-Find (o DSU) es una estructura de datos que mantiene un 
conjunto de elementos divididos en varios conjuntos disjuntos. 
Soporta dos operaciones principales:

1. Find: Determina a qué conjunto pertenece un elemento (devuelve
   el "representante" o "líder" del conjunto).
2. Union: Une dos conjuntos en uno solo.

Con optimizaciones como "Path Compression" y "Union by Rank", estas
operaciones se vuelven prácticamente instantáneas O(α(n)).
}

\diagrama{\begin{verbatim}
Estado Inicial: {0} {1} {2} {3} {4}
(Cada uno apunta a sí mismo como líder)

Union(0, 1): {0, 1} {2} {3} {4}
  [0] ← [1] (0 es el líder de 1)

Union(1, 2): {0, 1, 2} {3} {4}
  [0] ← [1] ← [2]

Path Compression (al buscar el líder de 2):
  [0] ← [1], [0] ← [2] (Ahora 2 apunta directo al líder 0)
\end{verbatim}}

\begin{lstlisting}[language=C++]
1. Find(i): Sigue los punteros `parent` desde `i` hasta llegar al 
   nodo que es su propio padre. Ese es el representante.
   - Optimización (Path Compression): Mientras subes, haces que 
     cada nodo apunte directo al representante final.
2. Union(i, j): Encuentras los representantes de `i` y `j`. Si son
   diferentes, haces que uno apunte al otro.
   - Optimización (Union by Rank): Siempre haces que el árbol más 
     bajito cuelgue del árbol más alto para mantenerlo balanceado.

📝 PSEUDOCÓDIGO
──────────────────────────────────────────────────────────────

// Find con Path Compression
función buscar(i):
    SI parent[i] == i: return i
    parent[i] = buscar(parent[i]) // ¡Compresión aquí!
    return parent[i]

// Union por Rango
función unir(i, j):
    raiz_i = buscar(i)
    raiz_j = buscar(j)
    SI raiz_i != raiz_j:
        SI rango[raiz_i] < rango[raiz_j]:
            parent[raiz_i] = raiz_j
        SINO:
            parent[raiz_j] = raiz_i
            SI rango[raiz_i] == rango[raiz_j]:
                rango[raiz_i]++

💻 CÓDIGO C++
──────────────────────────────────────────────────────────────

#include <iostream>
#include <vector>
using namespace std;

class UnionFind {
    vector<int> parent;
    vector<int> rank;
public:
    UnionFind(int n) {
        parent.resize(n);
        rank.resize(n, 0);
        for(int i = 0; i < n; i++) parent[i] = i;
    }

    int find(int i) {
        if (parent[i] == i) return i;
        return parent[i] = find(parent[i]); // Path compression
    }

    void unite(int i, int j) {
        int root_i = find(i);
        int root_j = find(j);
        if (root_i != root_j) {
            if (rank[root_i] < rank[root_j]) parent[root_i] = root_j;
            else if (rank[root_i] > rank[root_j]) parent[root_j] = root_i;
            else {
                parent[root_i] = root_j;
                rank[root_j]++;
            }
        }
    }

    bool connected(int i, int j) {
        return find(i) == find(j);
    }
};

int main() {
    UnionFind uf(5);
    uf.unite(0, 1);
    uf.unite(2, 3);
    
    cout << "¿0 y 1 conectados? " << uf.connected(0, 1) << endl; // 1 (true)
    cout << "¿0 y 2 conectados? " << uf.connected(0, 2) << endl; // 0 (false)
    
    uf.unite(1, 2);
    cout << "¿0 y 2 conectados ahora? " << uf.connected(0, 2) << endl; // 1 (true)

    return 0;
}

⏱️ COMPLEJIDAD (Big-O)
──────────────────────────────────────────────────────────────

  Operación   │ Tiempo             │ En la práctica
  ────────────┼────────────────────┼──────────────────────────
  Find        │ O(α(n))            │ Casi O(1)
  Union       │ O(α(n))            │ Casi O(1)
  Connected   │ O(α(n))            │ Casi O(1)
  Espacio     │ O(n)               │ Para guardar parents/ranks

* α(n) es la función inversa de Ackermann, que crece tan lento que 
  para cualquier valor real de `n`, es menor a 5.
\end{lstlisting}

\entrevistas{
✅ Cuándo usar: Problemas de conectividad, componentes conectados 
   en grafos, o detectar ciclos en grafos no dirigidos.

💡 Alternativa a DFS/BFS: Para saber si dos nodos están conectados,
   Union-Find es a veces más rápido y fácil de implementar que 
   un recorrido completo.

⚠️ Optimización: NUNCA olvides la "Path Compression". Sin ella, 
   Union-Find puede degradarse a O(n).

🔑 Problema Clásico: "Number of Islands" o "Redundant Connection".
}

\resumen{
• Gestiona grupos de elementos y sus uniones.
• Operaciones: `find` (quién es el jefe) y `union` (hacerse amigos).
• Casi O(1) gracias a Path Compression y Union by Rank.
• Ideal para detectar ciclos y conectividad en grafos.
• Solo funciona para conjuntos disjuntos (un elemento no puede estar
  en dos grupos a la vez).
════════════════════════════════════════════════════════════════
}

\end{document}